\documentclass[12pt,a4paper]{article}
\usepackage[scheme=plain]{ctex}
\usepackage[margin=2.5cm]{geometry}
\usepackage{indentfirst}
\usepackage{setspace}
\usepackage{fontspec}
\setCJKmainfont[BoldFont=SimHei,ItalicFont=KaiTi]{SimSun}
\newCJKfontfamily\kaifont{KaiTi}[AutoFakeBold]
\usepackage{newpxtext,newpxmath}

\usepackage{lastpage}





\pagestyle{plain}

\usepackage[nobottomtitles*]{titlesec}
\titleformat{\section}{\Large\bfseries}{}{0em}{}
\titlespacing*{\section}{0pt}{3ex}{1.5ex}

\usepackage{xcolor}
\usepackage{needspace}

\setlength{\parindent}{2em}
\setlength{\parskip}{0.4em}

\doublespacing

\title{\bfseries Day 1\\[0.3em]准备！踏足于权力之野}
\author{}
\date{}

\begin{document}
\maketitle

\vspace{-2em}

\vspace{1em}

\section{前哨的第一封信}
\vspace{-0.3em}{\color{gray}\small\kaifont\noindent 田野调查，即通过勘测、询问、交谈、观察等手段开展实地考察、获取研究资料的研究方法。K 是我们这支田野调查队的先遣队员，已经提前到达了"城堡"所在的村庄。这是他发回来的第一封信。他的任务是尽可能地感受这个地方、理解这个地方，然后把他看到的、想到的、困惑的东西传回来，帮助我们这些即将出发的人做好心理上和认知上的准备。读完这封信后问问自己：K 在向我们传递怎样的思想姿态？}\vspace{0.5em}
\vspace{0.3em}

\noindent\rule{\textwidth}{0.4pt}

\vspace{0.3em}



同志们：



我到了。时间很少，说关键的。



城堡看不见。不是被什么挡住了，是它本身就不让你从远处辨认。但你往下走，走进村子，会发现这里每一件小事都指向它。老板看你时的眼神，人们沉默的方式，一个年轻人半夜把我叫醒盘问时那种既傲慢又紧张的样子——他在维护某种他自己也说不清楚的东西。还有一通夜里打到总办公厅的电话，有人接。这些就是我的全部线索。



我开始意识到一件事。这里最重要的东西不直观。它在语气里，在沉默的长度里，在人们视线交错的方式里。你得让自己完全沉浸进去感受，又随时能抽身出来看清楚自己感受到了什么。这件事比我预想的难得多。



K

\newpage
\section{卢卡奇论卡夫卡}
\vspace{-0.3em}{\color{gray}\small\kaifont\noindent 卢卡奇是 20 世纪最重要的马克思主义文学理论家。在这段选文中，他通过讨论一个看似技术性的问题，小说的细节描写，揭示了卡夫卡对现实的特殊态度，即在虚无主义的基底上忠实反映扭曲的现实。我们将要采取的田野调查姿态与卢卡奇有何共同之处？又是否要保留一些必要的差异？}\vspace{0.5em}
\vspace{0.3em}

\noindent\rule{\textwidth}{0.4pt}

\vspace{0.3em}



同样的区分也适用于描写性细节。孤立地看，描写性细节可以是现实的足够真实的反映——前提是这位作家确实有才华。但描写的顺序和结构能否构成对客观现实的恰当呈现，则取决于这位作家对现实整体的态度。因为正是这种态度决定了每个个别细节在整体结构中被赋予什么样的功能。如果处理起来缺乏批判意识，其结果可能是一种任意的自然主义，因为作家将无法区分有意义的细节和无意义的细节。我认为乔伊斯就是一个典型的例子。这再次凸显了现代主义本质上自然主义的一面。



卡夫卡的情况更加复杂。卡夫卡是极少数在细节处理上具有选择性的现代主义作家之一，而不是自然主义的。从形式上看，他的细节处理与现实主义作家并无太大不同。只有当检视他的基本承诺——那条决定细节选择与编排顺序的原则——时，差异才会浮现。对卡夫卡而言，这条原则是他对超越性力量（虚无）的信念；因此，在他那些虚无主义的寓言中，艺术的统一性被打破了。



但这个问题不能以形式主义的方式去处理。在有些伟大的现实主义作品中，直接的社会历史现实被超越，细节上的现实主义是建立在对超自然世界的信念之上的。以 E.T.A. 霍夫曼为例。在霍夫曼那里，细节上的现实主义与对现实之幽灵本质的信念并行不悖。不过，细察之下，他的艺术目标与现代主义之间的差异便显现了。霍夫曼的世界——尽管带着童话般的、幽灵化的氛围——相当准确地反映了他那个时代德国的状况：一个正从扭曲的封建绝对主义走向同样扭曲的资本主义的国家。在霍夫曼那里，超自然是一种呈现德国整体状况的手段——在一个社会条件尚不允许直接现实主义描写、甚至尚未揭示出典型模式的时代。典型化的锻造在更为发达的法国要容易得多——尽管巴尔扎克偶尔也使用霍夫曼发展出来的方法（如在《和解的梅尔莫特》中）。



卡夫卡比霍夫曼更世俗。他的幽灵属于日常的资产阶级生活；而既然这种生活本身就已经是不真实的，便不需要霍夫曼式的超自然幽灵。但世界的统一性被打碎了，因为一种本质上主观的视像被等同于现实本身。帝国资本主义世界（其后来法西斯产物的先兆）所制造的恐怖——人被降格为纯粹对象——这种恐惧，原本是一种主观体验，却成为了一个客观实体。对一种扭曲的反映，变成了扭曲的反映。尽管卡夫卡的艺术方法与其他现代主义作家有所不同，但其呈现的原则是相同的：世界是超越性虚无的寓言。在卡夫卡的后继者那里，这些差异逐渐缩小或完全消失。以贝克特为例，他混合了卡夫卡式的和乔伊斯式的母题，最终产品是一种完全标准化的虚无主义现代主义。
\par\vspace{0.3em}{\footnotesize\color{gray}（Lukács, G., \textit{The Meaning of Contemporary Realism}, Merlin Press, 1963。由 DeepSeek V4 Pro 翻译。）}
\newpage
\section{观察与干预：田野中的伦理困境}
\vspace{-0.3em}{\color{gray}\small\kaifont\noindent 这段来自田野调查圣经的选文讨论了一个看似微小但至关重要的问题：研究者在田野里掏出笔记本的那一刻，发生了什么？她与被观察者的关系在那一瞬间变了——信任可能破裂，角色可能重构。然而，无涉的观察可能是不存在的，在充斥策略与抉择的田野调查中尤为如此...}\vspace{0.5em}
\vspace{0.3em}

\noindent\rule{\textwidth}{0.4pt}

\vspace{0.3em}



记速记不是简单地在笔记本或电脑上写几个字的事。由于速记往往是在研究对象附近、甚至就在他们说话做事的当下写的，因此写速记本身就是一个社会性和互动性的过程。具体而言，研究者如何记、何时记速记，可能对那些人如何看待和理解她是谁、她在做什么，产生重要的影响。关于是否该记速记、如果记的话何时记怎么记，并没有什么硬性规则。但随着在某个场景中待的时间增长，通过不断试错的好处，田野研究者可以逐渐演化出一套独特的做法，使写速记这件事适应那个场景的轮廓和限制。



\vspace{0.2em}



田野研究者还必须决定何时、何地、如何写速记。显然，低头看本子或键盘去写速记会分散研究者的注意力（即便只是一瞬间），使得对他人可能正在进行的复杂、快速而微妙的行动进行密切、持续的观察变得非常困难。但除了注意力的局限之外，记速记的决定还会对与田野中人们的关系产生重大影响。研究者努力与参与者建立紧密的联系，以便能被纳入他们生活中核心的活动。然而，在这些活动进行的过程当中，她可能体验到深深的矛盾：一方面，她可能希望通过趁话语刚说出口就记下来、趁场景正在上演就描摹下来，来保存当下的即时性；另一方面，她又可能觉得掏出一个笔记本或智能手机会毁掉这个瞬间，播下不信任的种子。参与者此时可能会把她看作一个主要兴趣在于发现他们的秘密、把他们最私密和最珍视的经历变成科学探究对象的人。



几乎所有民族志研究者在某些时刻都会感到被撕裂，一方面是自己的研究承诺，另一方面是自己真诚地融入所进入的那个世界的愿望。为了解决这些棘手的关系和道德问题，许多研究者主张：如果那些被研究的人没有完全知情和明确同意，就不应开展研究的任何环节。在这种观点看来，场景中的人必须被理解为合作者，他们与研究者积极合作，向外界讲述自己的生活和自己的文化。这种相互合作要求研究者请求许可才能撰写有关事件的内容，同时也要求尊重人们不愿揭示自己生活某些方面的意愿。



另一些田野研究者则觉得不必如此严格地被征求许可的要求所约束。一些人为这种立场辩护时坚持认为，田野研究者没有披露自己意图的特殊义务，因为所有的社会生活都包含掩饰的成分，没有人会完全暴露自己所有更深层的意图和私下的活动。另一些研究者指出，为自己而写的速记和田野笔记不会对他人造成直接伤害。这种做法当然是把那些艰难的道德和个人问题推迟到日后面对是否出版或以其他方式公开这些写作的决定时再去应对。最后，还有一些人主张向当地人隐瞒自己的研究目的，理由是所获得的信息将服务于更大的善。例如，如果研究者想要描述和公开无证工厂工人或养老院中老人的生存状况，他们就必须对控制这些场所准入权的当权者隐瞒自己的意图。
\par\vspace{0.3em}{\footnotesize\color{gray}（Emerson, Fretz \& Shaw, \textit{Writing Ethnographic Fieldnotes}, 1995。由 GLM 5.2 译。）}
\newpage
\section{田野调查程序清单}
\vspace{-0.3em}{\color{gray}\small\kaifont\noindent 《文化人类学调查——正确认识社会的方法》是国内多所高校人类学、社会学专业田野调查方法的常用参考书。以下是该书第二章"方法论（上）"的完整目录。那么，田野调查究竟是一个怎样的过程？这个过程在何种意义上是可规划或应该被规划的？}\vspace{0.5em}
\vspace{0.3em}

\noindent\rule{\textwidth}{0.4pt}

\vspace{0.3em}



\vspace{0.3em}\textbf{第二章 方法论(上)……………………………(14)}



\noindent\hspace{1em}一、调查课题的选择…………………………………(14)



\noindent\hspace{1em}二、调查方案和调查提纲的拟定……………………(15)



\noindent\hspace{1em}三、调查前的准备工作………………………………(17)



\noindent\hspace{2.5em}(一)了解地理和历史的背景材料…………………………(17)

\noindent\hspace{2.5em}(二)掌握前人调查研究成果………………………………(18)

\noindent\hspace{2.5em}(三)调查队伍的组成…………………………………………(18)

\noindent\hspace{2.5em}(四)装备、物品的购置……………………………………(19)



\noindent\hspace{1em}四、调查的开始…………………………………………(21)



\noindent\hspace{2.5em}(一)选点和安排生活………………………………………(21)

\noindent\hspace{2.5em}(二)说明来意………………………………………………(22)

\noindent\hspace{2.5em}(三)熟悉环境………………………………………………(22)

\noindent\hspace{2.5em}(四)查阅地方文献…………………………………………(23)

\noindent\hspace{2.5em}(五)人口调查………………………………………………(23)



\noindent\hspace{1em}五、如何接近人民及选择报告人……………………(24)



\noindent\hspace{1em}六、关于学习语言……………………………………(26)



\noindent\hspace{1em}七、关于参与观察……………………………………(27)



\noindent\hspace{1em}八、关于访问…………………………………………(28)



\noindent\hspace{1em}九、关于抽样…………………………………………(30)



\noindent\hspace{1em}十、材料的记录和核实………………………………(31)



\noindent\hspace{2.5em}(一)影视记录………………………………………………(32)

\noindent\hspace{2.5em}(二)声音记录………………………………………………(32)

\noindent\hspace{2.5em}(三)文字记录………………………………………………(33)

\noindent\hspace{3.5em}笔记　实录　原文　日记　图表　谱系

\noindent\hspace{2.5em}(四)关于核实材料…………………………………………(35)



\noindent\hspace{1em}十一、实物标本的搜集………………………………(35)



\noindent\hspace{1em}十二、调查成果的整理公布和存档…………………(38)
\par\vspace{0.3em}{\footnotesize\color{gray}（汪宁生，《文化人类学调查——正确认识社会的方法》，文物出版社，1996）}
\newpage
\section{K 的第一份纪实}
\vspace{-0.3em}{\color{gray}\small\kaifont\noindent K 深夜抵达村庄，被一个叫施瓦尔策的年轻人叫醒盘问，一通打到城堡总办公厅的电话确认了他的"土地测量员"身份。然而，K 真的是土地测量员吗？}\vspace{0.5em}
\vspace{0.3em}

\noindent\rule{\textwidth}{0.4pt}

\vspace{0.3em}



K到达时，已经入夜了。村子被厚厚的积雪覆盖着。城堡所在的山岗连影子也不见，浓雾和黑暗包围着它，也没有丝毫光亮让人能约略猜出那巨大城堡的方位。K久久伫立在从大路通往村子的木桥上，举目凝视着眼前似乎是空荡荡的一片夜色。

然后他去找过夜的地方；酒店里人们还都没有睡，店老板虽然无房出租，但在对这位晚客的突然到来感到极度惊讶和惶乱之余，还是愿意让他在店堂里垫一个装稻草的口袋睡觉，K同意这一安排。有几个农民还在喝啤酒，但他无意同任何人交谈，便自己去阁楼上把草袋搬下来，在炉子附近躺下了。屋里很暖和，那几个农民说话声音很低，他用困倦的双眼打量了他们一番便倒头睡下了。

然而没有多久他便被吵醒了。这时只见一个城里人装束、长着一副演员般的面孔、浓眉细眼的年轻人同店老板一起站在他身边。那些农民也都还没有走，其中几个把椅子转过来对着他们，以便看得更清楚、听得更真切些。年轻人为吵醒了K而十分客气地向他道歉，自我介绍说，他是城堡主事的儿子，然后说道：“这村子是城堡的产业，凡是在这里居住或过夜的，从某种意义上讲也算是在城堡居住或过夜，没有伯爵大人的许可，谁也不能在此居留。可是您并没有得到这样的许可，至少您并没有出示这样的证明嘛。”

K半坐起身子，捋了捋头发，仰头看着众人说道：“我迷了路，这是摸到哪个村子来了？这里是有一座城堡吗？”

“那还用问，”年轻人慢条斯理地说，这时店堂各处都有人大惑不解地冲着K摇头。“这里是伯爵大人威斯特威斯的城堡。”

“一定要得到许可才能在这儿过夜吗？”K问道，似乎想弄清刚才他听到的那些话是不是在做梦。

答话是：“是必须得到许可才行，”紧接着这个年轻人伸出胳臂，向店老板和酒客们问道：“难道竟有什么人可以不必得到许可吗？”那话音和神态里，包含着对K的强烈嘲笑。

“那么我只好现在去讨要许可了。”K打着哈欠说，一面推开被子，似乎想站起来。

“向谁去讨要？”年轻人问。

“向伯爵大人，”K答道，“恐怕没有什么别的法子了吧。”

“现在，半夜三更去向伯爵大人讨要许可？”年轻人叫道，后退了一步。

“这不行吗？”K神色泰然地说，“那么您为什么叫醒我？”

这时年轻人憋不住火了。“真是活脱脱一副盲流口吻！”他喊叫起来，“伯爵衙门的尊严必须维护！我叫醒您是想告诉您：您必须立即离开伯爵领地！”

“好了，戏做够了吧。”K用异常轻的声音说，接着又躺下去，拉过被子盖在身上。“年轻人，您太过分了点，我明天还要再考虑考虑您今天的表现的。如果一定要见证的话，酒店老板和这里的各位先生就可以作证。现在请您听清楚：我是伯爵招聘来的土地测量员，明天我的助手就要带着仪器乘车随后跟来。我因为不想失去这个踏雪觅途的好机会，所以步行前来，可惜几次迷路，才到得这样晚。现在到城堡去报到时间已经太迟，这一点我自己很清楚，用不着您来赐教。正因为这样我才勉强在这草袋上凑合过夜，而您竟然——客气点说吧——举止失礼，打搅我休息。好了，我的话说完了，晚安，诸位先生！”说到这里K翻了一个身，转向炉子去了。

“土地测量员？”他听见背后将信将疑地发问，过后又无人做声了。然而不久那位年轻人便克制住自己，用压低了的——低到可以被看作是为了照顾K的睡眠，然而又大到能让他听清楚——的声音对老板说道：“我现在就打电话去问一下。”怎么，在这个乡村小酒店里居然还有电话？唔，设备还真够齐全的。这些事，一件一件地听来也使K感到惊奇，不过总起来却又在他的意料之中。他发现，电话差不多正好摆在他的头顶上方，刚才他睡眼惺忪，没有注意到。现在，如果年轻人一定要打电话，那么他无论如何不能不打搅K的睡眠，问题只在于K让不让他打这个电话？K决定还是让他打。可是这样一来，在底下装睡便没有什么意思了，于是K又恢复了仰卧的姿势。他看见农民们怯怯地聚到一起，嘁嘁喳喳议论，看来土地测量员的到来不是一件小事。这时厨房的门早已打开，膀大腰圆的老板娘站在那里几乎把门框塞满，店老板踮着脚尖走到她跟前去告诉她刚才发生的一切。现在，电话接通，开始通话了，城堡主事已经就寝，是一位副主事——数位副主事当中某位名叫弗里茨的老爷——在那边接电话。年轻人自报姓名，说是叫施瓦尔策，接着便说他发现有一个名叫K的三十多岁的男子，衣冠不整，心安理得地在酒店里一个草袋上睡觉，枕着一个小得可怜的背囊，手边放着一根拐杖。他施瓦尔策自然觉得此人形迹可疑，而因为店主在这件事上显然失职，他施瓦尔策当然就责无旁贷地要过问此事，查明情况了。对于被叫醒盘问，对于他施瓦尔策按职责惯例作出的要将K逐出伯爵领地的警告，K的反应是很不耐烦，总之看起来火气不小，然而也许不无道理，因为他自称是伯爵大人聘来的土地测量员。在这种情况下，当然至少从例行手续上看有必要对他的话进行核实，因此，他施瓦尔策恳请弗里茨老爷在城堡总办公厅询问一下是否确有这样一位土地测量员应聘前来，并请将答复立即电话告知。

电话打完店堂里便安静下来，弗里茨在那边询问，这边在等着回话。K的神态一如既往，连头也不回，看样子一点也不急于知道结果如何，两眼茫然直视前方。施瓦尔策对事情的描述是一种不怀好意和谨小慎微的混合物，这番话给了K一种印象：城堡里连施瓦尔策这样的小人物也能很容易得到可说是外交手腕方面的训练。另外，看起来那里的人也决不偷懒；你看，总办公厅有人上夜班呢。而且显然很快就做出回答，因为这时弗里茨打电话来了。不过，好像这个回话太简短，因为施瓦尔策马上又气呼呼地把听筒挂上。“我早就说了嘛！”他大声叫道：“土地测量员，连影子也没有！这人真是个卑鄙的、信口雌黄的流浪汉，说不定还更坏呢！”此时K脑子里闪过一个念头，他想：所有的人，就是施瓦尔策、那些农民，还有店老板、老板娘，眼看就要向他猛扑过来。为了至少避一避这个凶猛的势头，他从头到脚，整个儿蜷缩到被窝里面去了。这时电话铃又响起来，而且，K觉得声音特别急促响亮。于是他又把头从被子里慢慢伸出来。虽然电话几乎不可能仍是涉及K的事情，但所有的人还是一下子突然肃静下来，施瓦尔策则回到电话机旁去。他站在那儿耐心地听完了一番较长的解释之后，低声说道：“那么是弄错了？我实在是太难为情。办公室主任亲自打来了电话？真是怪事，真是怪事。我该怎么向土地测量员先生解释才好呢？”

K听了这些话精神为之一振。这么说，城堡已经任命他为土地测量员了。这个情况一方面对他不利，因为它表明，城堡已经了解他的底细，在反复掂量了双方的力量对比之后，决定成竹在胸地开始同他较量。可是另一方面，情况又对他有利，因为在他看来这同时也证明对方低估了他，他将会有兴许比他自己起初敢于希冀的还要更多的一些自由。并且，如果以为通过这从心理战来看不能不说是高明的一着，即公开承认他的土地测量员身份，就能使他经常处于诚惶诚恐、心惊胆战的状态的话，那他们就错了；这一官方认可使他吃了一惊，但也就止此而已。

对羞怯地向他走来的施瓦尔策，K摆手示意他不必过来了；此刻人们忙不迭地恳求他赶快搬到老板的房里去住，他也拒绝了，只是接受了老板递给他的一杯催眠饮料，又从老板娘手里接过一盆水、肥皂和毛巾。现在根本用不着他开口叫人离开店堂，原来这时所有的人都把脸背过去使劲往外挤，可能是怕明天被他认出来吧。关灯之后，他终于可以休息了。他睡得很香，除了一两次被跑过的老鼠惊醒之外，一直睡到第二天早晨。
\par\vspace{0.3em}{\footnotesize\color{gray}（第一章，1-5页）}
\newpage
\section{官僚制的混乱性（Kafkaesque）}
\vspace{-0.3em}{\color{gray}\small\kaifont\noindent 这是一个印度农民的“上访”故事。研究者将这个故事归结于卡夫卡式的官僚机构的混乱运作。你认可吗？}\vspace{0.5em}
\vspace{0.3em}

\noindent\rule{\textwidth}{0.4pt}

\vspace{0.3em}



2013至2014年冬季，Shreegarh 村 15 户人家的农田被政府征用，事先没有任何通知。被征土地面积因户而异，总共约 11 英亩，户均不到 1 英亩。政府征地的目的是修路。受影响的农民均来自低种姓，主要是 kushwaha 和 vishwakarma（badhai）种姓。冬季是本德尔坎德地区的小麦种植季节，也是农民的主要收入来源。土地被征走时，小麦正在地里生长。过去曾在类似的国家干预下被迫搬迁的经历，让整个村庄弥漫着一种“似曾相识”的感觉。正如卡夫卡《审判》的主人公约瑟夫·K——“他什么也没有做错，然而，一天早晨，他被逮捕了”——这 15 户村民的土地在不经通知、不作补偿的情况下就被征走了。



一位名叫 Ramraja 的低种姓农民告诉我们，村民之间有一种普遍的看法：军队会帮他们，因为上世纪 60 年代他们村子的土地就是被政府征用去建军事设施的。村民们并不清楚应该找哪个政府部门来解决问题。30 岁的低种姓农民 Murari 说，曾有一支军队驻扎单位在村里举办过活动，当时宣布要“认领”这个村庄，当地报纸还报道了这件事。不知道该怎么办的村民们，拿着那份剪报来到了军营驻地。Murari 和另外两个村民一起去见了一位军官。他告诉我们：



\begin{quotation}\kaifont\noindent 我们给他（军官）看了那份剪报。军官让我们去法院。他说："我们收的是 bhumidhar 地（可在农民名下登记的土地），还的也是 bhumidhar 地。现在如果你们或政府把它作为安置地来保留，那就不在我们的管辖范围之内了。"我们没法跟 sahib log（长官们）争辩。我们也去找过一个政治领袖，但没有人帮忙。\end{quotation}



bhumidhar 地的所有者拥有出售土地的法定权利。但如前所述，分配给 Shreegarh 村民的土地不可买卖。邻近村庄（如 Dosah）的地价在上涨，Shreegarh 的村民却无法从中受益。在底层语境中，质疑地方精英和政府官员被视为一种冒犯。就像《城堡》中那些村民一样，Shreegarh 的底层民众在与官员争论时，越过某一界限便会退缩。他们曾前往地区行政长官（DM）办公室要求补偿。村民们告诉我们，主管该地区行政事务的 DM 拒绝对此承担任何责任，理由是征地“依法进行”。他还说，如果村民要继续追究，可以去三百多公里外的高等法院起诉。



村民们聘请了一位来自附近村庄、在 Rajapur 执业的律师。律师向他们收取了近 10 万卢比，但案子至今连高等法院的门都没进。Murari 说，律师向村民们保证案件几个月内就能解决，并要求每个受影响的农民缴纳 3000 卢比，之后又不断索要更多。受影响的村民最近见到了这位律师，令他们沮丧的是，律师竟然要求他们重新提供土地记录的复印件和身份证件——换句话说，律师连在高等法院立案的程序都还没启动。



印度的司法系统被积压案件严重压垮，从指控确定到审判开始之间隔着极为漫长的时间。印度报业托拉斯 2014 年的报道指出，印度有 3000 万件积压案件，其中 2600 万件积压在基层法院。Kasturi（2009）估计民事案件积压时间从 9 个月到 5.4 年不等。案件的积压也导致了庭外解决倾向的增加。律师的拖延和不作为对村民们来说不是一个好兆头——他们获得正义的希望渺茫。对于 Shreegarh 的农民来说，法律就像一本字迹消褪又覆盖的羊皮卷，如同《在流放地》中那位指挥官无法辨认的图纸。协议是用英文写成的——而且是法律英文——难以理解。村里没有人精通英语去破译协议上的文字。法律语言和法律程序对村民们而言始终不可辨认。当判决以国家征地的具体形式施加到他们身上时，他们对法律知之甚少。



法律的不可接近性还体现在村民寻求帮助的各种失败尝试中。各级守门人挡住了他们的去路。地区行政部门、Lekhpal（土地记录官）、政府官员构成了阻挡村民接近正义的第一层守门人。就像卡夫卡《在法的门前》中那个守门人一样，地区行政部门把门留着——说村民们可以去法院起诉。但与此同时，土地被征了，路已经修好了，农民的损失已成既定事实。还有更强大、更高的守门人守护着法律——法院。印度的司法体系由多个层级组成。最高法院位于首都新德里。最高法院之下是各邦的高等法院。地区法院处于下一级，设在各个县府所在地。地区法院下还有审理民事和刑事案件的次级法院。高等法院——按照那位律师的说法，Shreegarh 农民的案件应该在这里提交——距离村子三百多公里。除了长途跋涉，高等法院的听证又是极其昂贵的。受影响的家庭承受不起这种法律程序的全部成本。尽管法律无处不在，对村民而言却始终遥不可及。



在卡夫卡的世界里，法律采取了自然法的形式，没有错误。国家官僚机构从同一种法律中汲取权威，也绝不承认错误。这种理性在这里以一种循环的形式展开：首先，官僚机构拒不承认自身有任何过失；其次，即便存在某种不规范，纠正的程序本身就是一个法律问题，而这个问题对穷人来说几乎不可接近。因此，官僚机构就像《城堡》里那样——变成了自我服务的机器。结果是规则变得不透明，法律对底层民众而言不可接近，就像那个在“法的门前”等待了一生的乡下人。复杂的程序和晦暗的规则，使法律和官僚程序对底层民众来说始终不可辨认。
\par\vspace{0.3em}{\footnotesize\color{gray}（Khare, A., \& Varman, R., "Kafkaesque institutions at the base of the pyramid", \textit{Journal of Marketing Management}, 2016。由 DeepSeek V4 Pro 翻译。）}
\end{document}